\section{The experiments}
\label{sec:experiments}

The experiments of the paper are programs, not hand-made runs: anyone can
run them again and get the same files. This section explains, for each
one, the question it answers, what varies, what stays fixed, what it
measures and how to read its output.

\subsection{How an experiment runs}
\label{sec:ex-how}

\begin{figure}[ht]
\centering
\begin{tikzpicture}[node distance=4mm]
\node[box,minimum width=32mm,text width=30mm] (pts) {\textbf{design points}\\\scriptsize one parameter setting each (a label, a rule)};
\node[box,minimum width=32mm,text width=30mm,below=6mm of pts] (inst) {\textbf{instances}\\\scriptsize applications generated from (seed, $i$)};
\node[box,right=14mm of $(pts)!0.5!(inst)$,minimum width=34mm,minimum height=22mm,
	fill=brand,text=white,text width=32mm] (run) {\textbf{every point $\times$ every instance}\\\scriptsize in parallel, one design point per thread};
\node[file,right=12mm of run,yshift=7mm,minimum width=34mm] (csv) {\textbf{results/NAME.csv}\\\scriptsize one row per design point};
\node[file,right=12mm of run,yshift=-7mm,minimum width=34mm] (tab) {\textbf{summary table}\\\scriptsize means per point};
\draw[arrow] (pts.east) -- (run.west |- pts.east);
\draw[arrow] (inst.east) -- (run.west |- inst.east);
\draw[arrow] (run.east |- csv.west) -- (csv.west);
\draw[arrow] (run.east |- tab.west) -- (tab.west);
\end{tikzpicture}
\caption{An experiment: every design point is run on the same generated
applications, so that two points differ only by their parameters.}
\label{fig:experiment}
\end{figure}

\begin{shell}
build/bin/experiment NAME [-{}-instances N] [-{}-jobs N] [-{}-seed N] [-{}-out FILE] [-{}-quiet]\\
build/bin/experiment all \quad\# every experiment, about three minutes on two cores
\end{shell}

An experiment is a list of \emph{design points}; each is run on several
generated applications (the \emph{instances}, 20 by default). Instance
$i$ always uses the application generated from (experiment seed, $i$):
every point sees the same applications, and the results do not depend on
the number of threads. Unless an experiment says otherwise, a point uses
the defaults: 300 task-chains of the default generator, the built-in
zonal platform (8 big cores on the PCU, 4 microcontroller cores at 50\,\%
on each of the five zone controllers), $\varphi = 0.8$, $\sigma = 0.4$,
LET tasks of 50\us, a star at 1\,Gbit/s with one CBS class, and every
step of the method.

The program prints a progress line, then a table with the means of the
metrics of the experiment over the instances of each point (over the
feasible ones for the mapping metrics), and writes every design point to
\cmd{results/NAME.csv} (\cref{sec:ex-csv}). The columns of the tables:

\begin{center}\small\row
\begin{tabularx}{\textwidth}{@{}>{\raggedright\arraybackslash}p{36mm}X@{}}
\toprule\rowcolor{white}
column & meaning (mean over the instances) \\
\midrule
feasible \% & instances where every sub-chain found a core \\
cores $M$, lower bound & cores used, and the lower bound on them \\
PCU, ZCU cores & cores used on the central computers, on the other ECUs \\
FR \% & fragmentation rate: cuts in \% of the possible cuts \\
SR sched, SR \% & schedulable cores after the priority assignment, and at the end \\
net hops, port max \% & inter-ECU communications; LET load of the busiest port \\
on time, misses & task-chains meeting their deadline (\%), late ones \\
reaction max & the largest maximum reaction time (ms) \\
scaling & critical scaling factor \\
CPU s & CPU time of a design point \\
\bottomrule
\end{tabularx}
\end{center}

\subsection{Experiments of the method}

\begin{experimentbox}{sweep}{the sub-chain limit against the rule}
\exitem{Question.} How does the limit $\sigma$ trade cores against
communications, and which rule packs best?
\exitem{Varies.} $\sigma \in \{0.02, 0.03, 0.05, 0.1, 0.2, 0.4, 0.8\}$
$\times$ the four rules: 28 points.
\exitem{Measures.} Cores and lower bound, cores per ECU kind, fragmentation
rate, schedulable cores, chains on time, maximum reaction time, CPU time.
\exitem{Reading.} A small $\sigma$ cuts chains into many small sub-chains:
fewer cores (they pack like sand), but more communications and longer
reaction times. Most chains use less than 8\,\% of a core, so the cuts
happen below $\sigma \approx 0.1$. WFDU uses every core of the platform
(28): it spreads on purpose.
\end{experimentbox}

\begin{experimentbox}{scaling}{CPU time against the size}
\exitem{Question.} Is the method fast enough for a whole vehicle?
\exitem{Varies.} 50 to 3200 task-chains (doubling) $\times$ the four rules;
2 instances per point.
\exitem{Measures.} Chains, tasks, feasibility, cores, schedulable cores,
CPU time of the placement phases and of the whole design point.
\exitem{Reading.} The CPU time grows a little faster than linearly: about
2\,s for 2000 chains on one core. The exact solver needs seconds to a
minute for 300 chains (\cref{sec:ex-exact}).
\end{experimentbox}

\begin{experimentbox}{architecture}{speed of the zone controllers}
\exitem{Question.} How fast must the zone controllers be?
\exitem{Varies.} Zone core speed 25, 50, 75, 100\,\% (PCU at 100\,\%)
$\times$ the four rules.
\exitem{Measures.} Cores (total, PCU, zones), lower bound, schedulable
cores, chains on time, LET load.
\exitem{Reading.} Slow zone cores take more cores for the tasks bound to
the zones; the free sub-chains go to the PCU anyway.
\end{experimentbox}

\begin{experimentbox}{affinity}{allocation with and without affinity}
\exitem{Question.} Does trying the cores of the other sub-chains of the
same chain first save communications?
\exitem{Varies.} $\sigma \in \{0.2, 0.4, 0.6\}$ $\times$ affinity on/off
$\times$ the four rules.
\exitem{Measures.} Cores, communications inside an ECU and over the
network, schedulable cores, chains on time, reaction time.
\exitem{Reading.} Compare each point with its twin without affinity: the
same cores, fewer hops.
\end{experimentbox}

\begin{experimentbox}{ablation}{each step of the method removed}
\exitem{Question.} What does each step of \cref{fig:method} bring?
\exitem{Varies.} Seven variants: the full method; without affinity;
without consolidation; with the optional swap moves; without relocation moves;
without the local search; without any of them (bin packing and the
priority assignment only), on two settings, \emph{default} and
\emph{stressed} (cores filled up to $\varphi = 1$, deadlines at 200\,\% of
the pattern periods instead of 300\,\%), $\times$ FFDU and CADU: 28
points.
\exitem{Measures.} Cores, schedulable cores after the priorities and at
the end, network hops, chains on time and late, CPU time.
\exitem{Reading.} Compare each variant with \emph{full}. With the default
setting, the consolidation saves cores (14.4 instead of 15.8 for FFDU) and
the local search brings chains on time (92.6 instead of 88.4\,\%). The
relocation moves only matter when cores are unschedulable, which happens
in the stressed setting (without them, 0.7\,\% of the cores stay
unschedulable). The swap moves change nothing, which is why they are off
by default.
\end{experimentbox}

\begin{experimentbox}{practice}{the method against practitioner rules}
\exitem{Question.} Does the method beat what an engineer would do by
hand?
\exitem{Varies.} The method with FFDU and with CADU; \emph{whole chains}
(cut only where the ECUs force it: $\sigma = \varphi$); \emph{central
only} (the free software on the central computer, the zones only for the
tasks bound there); both together; \emph{load balancing} (WFDU).
\exitem{Measures.} Cores (total, PCU, zones), lower bound, hops, busiest
port, chains on time, reaction time.
\exitem{Reading.} Keeping chains whole costs little here (the chains are
small), but see \cmd{limit}; putting all free software on the central
computer costs cores, network traffic and deadlines; load balancing uses
every core.
\end{experimentbox}

\begin{experimentbox}{limit}{the sub-chain limit where it matters}
\exitem{Question.} Why cut task-chains at all, and is the default
$\sigma = 0.4$ right?
\exitem{Varies.} $\sigma \in \{0.4, 0.6, 0.8\}$ (0.8 is the cap: whole
chains) on three settings: the default point, shared-memory contention at
50\,ns, and 600 task-chains; $\times$ FFDU and CADU: 18 points.
\exitem{Measures.} Feasibility, cores and lower bound, hops, chains on
time, CPU time.
\exitem{Reading.} At the default point a larger $\sigma$ changes little
(CADU saves about a core). With contention, the WCETs grow once the
sub-chains are placed, and a large sub-chain no longer fits any core:
only 45\,\% of the instances are feasible at $\sigma = 0.8$, 95\,\% at
0.6, all at 0.4. That is why the default is 0.4.
\end{experimentbox}

\begin{experimentbox}{deadlines}{deadlines and the local search}
\exitem{Question.} How tight can the deadlines be, and how much does the
local search recover?
\exitem{Varies.} Deadlines at 200, 300, 400, 600\,\% of the sum of the
pattern periods of each chain $\times$ local search on/off $\times$ the
four rules.
\exitem{Measures.} Cores, chains on time, late chains, moves kept, reaction
time, CPU time.
\end{experimentbox}

\begin{experimentbox}{fill}{the utilization cap}
\exitem{Question.} How high may the cores be filled?
\exitem{Varies.} $\varphi \in \{0.6, 0.7, 0.8, 0.9, 1.0\}$ $\times$ the four
rules.
\exitem{Measures.} Cores, schedulable cores after the priorities and at the
end, scaling factor, chains on time, CPU time of the placement.
\exitem{Reading.} Up to 0.9 the optimal priority assignment makes every
core schedulable; at 1.0 about 1\,\% of the cores are not (the periods are
not harmonic), and the relocation moves repair all of them.
\end{experimentbox}

\begin{experimentbox}{let}{the cost of the LET tasks}
\exitem{Varies.} WCET of the LET task: 10, 50, 100, 200\us{} $\times$ the
four rules.
\exitem{Measures.} Cores, LET load, communications on a core, schedulable
cores, chains on time.
\end{experimentbox}

\begin{experimentbox}{communication}{implicit communication or strict LET}
\exitem{Question.} What does the determinism of strict LET cost, against
implicit communication?
\exitem{Varies.} The two schemes on three settings: the default point,
deadlines at 500\,\%, and a ring with three classes; $\times$ FFDU and
CADU: 12 points.
\exitem{Measures.} Cores, hops, largest network bound, reaction times,
chains on time.
\exitem{Reading.} The cores do not change. Strict LET removes the response
time from the jitter of the flows, so the largest network bound falls by
15 to 19\,\%, but every crossing between cores waits for the end of the
producer's period: 30 points fewer chains on time with FFDU, 16 with CADU,
and only 4 and 1.5 with deadlines at 500\,\%.
\end{experimentbox}

\begin{experimentbox}{margins}{WCET margin or sizing}
\exitem{Question.} Two ways to keep a margin: raise every WCET at design
time, or size for a scaling factor. Which one costs fewer cores for the
same margin?
\exitem{Varies.} WCET margin 0, 10, 20, 40\,\%; then targets 1.10, 1.25,
1.50 for the sizing; $\times$ the four rules.
\exitem{Measures.} Cores, application load, schedulable cores, scaling
factor, cap found, target reached, chains on time.
\end{experimentbox}

\subsection{Experiments of the platform}

\begin{experimentbox}{cores}{types of cores}
\exitem{Question.} What do little cores and lock-step cores change?
\exitem{Varies.} Four platforms: PCU with 8 big cores or 4 big + 8 little;
zones with 4 microcontroller cores or 2 + 2 lock-step, the safety chains
(10\,\%) on lock-step cores only; every WCET up to 30\,\% off its speed
ratio on each type; $\times$ the four rules.
\exitem{Measures.} Cores (total, PCU, zones), application load, scaling
factor, chains on time.
\end{experimentbox}

\begin{experimentbox}{constraints}{spatial constraints}
\exitem{Varies.} None; 10, 20, 30\,\% of replicated chains (on other
ECUs); 20\,\% replicated on other cores only; 20\,\% of zone-only chains;
30\,\% of co-located tasks; 8\,MB of memory per zone controller; $\times$
the four rules.
\exitem{Measures.} Cores (total, PCU, zones), hops, chains on time,
reaction time.
\exitem{Reading.} Each constraint takes freedom from the allocation: its
price is read in cores and hops against the point \emph{none}.
\end{experimentbox}

\begin{experimentbox}{contention}{shared-memory contention}
\exitem{Varies.} Memory latency 0, 20, 50, 100\,ns per request (1000
requests per ms of WCET), $\times$ the four rules.
\exitem{Measures.} Cores (total, PCU, zones), application load, schedulable
cores, scaling factor, chains on time.
\exitem{Reading.} Contention grows with the cores used on an ECU, and more
load needs more cores: a packed mapping pays the most.
\end{experimentbox}

\begin{experimentbox}{dse}{four E/E architectures}
\exitem{Question.} For the same vehicle functions, which architecture
needs the fewest cores and the least network traffic, and up to what size
of software does each one work?
\exitem{Varies.} The four templates of \cref{fig:templates}
(distributed, domain, hybrid, zonal) $\times$ 100, 200, 300 task-chains
$\times$ FFDU and CADU: 24 points. The generator draws the same functions
for every template (the same seeds); their sensors and actuators are
wired to the ECUs that are not central computers.
\exitem{Measures.} Feasibility, cores and lower bound, hops, busiest port,
chains on time, reaction time, CPU time.
\exitem{Reading.} The distributed architecture runs out of cores as the
software grows (each small ECU has at most two cores, the ADAS controller
four): every instance is feasible at 100 chains, half of them with FFDU
at 300. The domain, hybrid and zonal ones all absorb 300 chains; the
hybrid and domain ones need about a core less than the zonal one, because
part of the bound work runs on big cores.
\end{experimentbox}

\begin{experimentbox}{lifecycle}{the SOP, then software updates}
\exitem{Question.} How many updates does a core bought at the SOP buy?
\exitem{Varies.} The cap at the SOP: 0.5, 0.6, 0.7, 0.8 $\times$ the four
rules. The platform is then fixed (the cores the SOP mapping uses), and
ten updates of 15 new task-chains each are mapped on it at $\varphi =
0.8$, without moving what is already there.
\exitem{Measures.} Cores of the platform, scaling factor, updates absorbed,
and after each update its chains, scaling factor and chains on time.
\end{experimentbox}

\subsection{Experiments of the network}

\begin{experimentbox}{network}{link rate, preemption, message sizes}
\exitem{Varies.} 100 or 1000\,Mbit/s $\times$ frame preemption on/off
$\times$ the default messages or sensor data (a third of large messages, up
to 64\,KB) $\times$ the four rules.
\exitem{Measures.} Cores, hops, largest network bound, busiest port, chains
on time, reaction time.
\end{experimentbox}

\begin{experimentbox}{topology}{star, tree and ring}
\exitem{Varies.} Star; tree and ring with a backbone at 1 or 10\,Gbit/s;
tree and ring with ATS; $\times$ the four rules.
\exitem{Measures.} As for \cmd{network}.
\exitem{Reading.} On a ring the bursts of the flows grow at every hop and
the dependencies between ports are cyclic; ATS reshapes every flow at every
switch and removes that growth.
\end{experimentbox}

\begin{experimentbox}{classes}{traffic classes on a whole vehicle}
\exitem{Varies.} 2000 task-chains, 16 PCU cores; one CBS class at
1\,Gbit/s; then with a 10\,Gbit/s backbone: 1, 2 or 3 CBS classes, 1 or 3
TAS classes, 1 or 3 ATS classes; $\times$ the four rules; 2 instances.
\exitem{Measures.} Cores, largest bound, busiest port, chains on time and
late, CPU time.
\end{experimentbox}

\begin{experimentbox}{frer}{frame replication on the ring}
\exitem{Varies.} CBS or ATS $\times$ 0, 10, 20, 30\,\% of replicated chains
$\times$ FRER on/off $\times$ the four rules, on a ring.
\exitem{Measures.} Cores, hops, replicated flows, largest bound, busiest
backbone link, chains on time.
\end{experimentbox}

\subsection{Case studies and assumptions}

\begin{experimentbox}{detection}{a small detection chain}
\exitem{Varies.} An illustrative ADAS chain of nine tasks (not taken from
a benchmark; identical cores), $\varphi = 1$ and $0.8$, $\sigma$ from 0.2 to
$\varphi$ by 0.1, $\times$ the four rules; one instance.
\exitem{Measures.} Cores, schedulable cores, communications on a core,
fragmentation rate, scaling factor.
\end{experimentbox}

\begin{experimentbox}{vehicle}{the software of a whole vehicle}
\exitem{Varies.} 2000 task-chains (about 8700 tasks), 16 PCU cores, at the
reference speeds or with PCU cores four times faster (zones at 100\,\%);
$\times$ the four rules; 2 instances.
\exitem{Measures.} Tasks, cores and lower bound, cores per ECU kind,
schedulable cores, chains on time, busiest port, CPU time.
\end{experimentbox}

\begin{experimentbox}{robustness}{the assumptions of the generator}
\exitem{Question.} Do the conclusions depend on how the applications are
generated?
\exitem{Varies.} One assumption at a time, below and above its default:
chains local to one zone (0 or 50\,\% instead of 25), chains with a
sensor and an actuator in a zone (30/25 or 90/80\,\% instead of 60/50),
chain utilizations (halved or doubled), heavy chains (0 or 5\,\% instead of
2), message sizes (mostly small, or 40\,\% large); $\times$ FFDU and
CADU: 22 points.
\exitem{Measures.} Cores and lower bound, hops, busiest port, chains on
time, reaction time.
\exitem{Reading.} The values move, the ranking does not: at every point
FFDU uses the fewest cores, and CADU the fewest network hops and the
largest share of chains on time.
\end{experimentbox}

\subsection{The CSV files}
\label{sec:ex-csv}

{\raggedright Every row starts with the design point and its parameters:
\cmd{experiment, point, label, application, instance, seed, sigma,
algorithm, phi, pcu\_speed, zcu\_speed, let\_wcet, wcet\_margin, affinity,
repair, link\_mbps, preemption, deadline\_pct, error, topology,
backbone\_mbps, classes, shaping, frer, pcu\_memory, zcu\_memory,
separation, replica\_pct, colocate\_pct, zone\_only\_pct, pcu\_little,
zcu\_lockstep, type\_variation, lockstep\_pct, target\_scaling,
memory\_latency, access\_density}. The \cmd{label} tells the points
apart when a parameter has no column of its own (the template of
\cmd{dse}, the variant of \cmd{ablation}). The results follow.\par}

\begin{center}\small\row
\begin{tabularx}{\textwidth}{@{}>{\raggedright\arraybackslash}p{0.43\textwidth}X@{}}
\toprule\rowcolor{white}
columns & meaning \\
\midrule
\cmd{chains, tasks, feasible} & size of the application; 1 if every sub-chain was placed \\
\cmd{fragments, cuts, forced\_cuts, max\_cuts, fr} & partitioning; fragmentation rate (\%) \\
\cmd{cores\_total, cores\_arch, cores\_used, cores\_lb} & cores at the end, initial cores, cores used $M$, lower bound \\
\cmd{cores\_pcu, cores\_zcu, cores\_big \dots\ cores\_lockstep} & $M$ on the central computers and the other ECUs, and by core type \\
\cmd{sr\_sched, sr\_swap, sr\_disp, sr\_final} & schedulable cores after priorities, swaps, relocations, at the end (\%) \\
\cmd{sr\_chains, sr\_on\_time, deadline\_misses} & task-chains schedulable and meeting their deadline (\%); late ones \\
\cmd{repairs, emptied} & moves kept by the local search; cores emptied by consolidation \\
\cmd{u\_app, u\_let, u\_unused} & application load, LET load, room below the cap (cores) \\
\cmd{hops\_intra, hops\_inter} & communications inside an ECU, between ECUs \\
\cmd{net\_overloaded, net\_avg, net\_max} & flows without a bound; mean and largest flow bound (\us) \\
\cmd{port\_max, backbone\_max, ports\_overloaded, frer\_flows} & busiest port and backbone link (\%); overloaded ports; replicated flows \\
\cmd{latency\_avg/max, delay\_avg/max, react\_avg/max} & end-to-end response time, end-to-end delay, maximum reaction time (\us) \\
\cmd{scaling, sized\_phi, target\_met, sizing\_runs, t\_sizing} & critical scaling factor; result of the sizing \\
\cmd{cores\_c$k$, u\_app\_c$k$, scaling\_c$k$, memory\_c$k$} & per ECU $k = 0..5$ (the first six ECUs of the platform) \\
\cmd{port\_out\_c$k$, port\_in\_c$k$} & load of the port of ECU $k$ and of the switch port towards it (\%) \\
\cmd{updates, update\_chains, update\_phi, absorbed, platform\_cores} & lifecycle: updates tried, chains per update, cap, updates absorbed, cores of the platform \\
\cmd{life\_scaling\_$u$, life\_on\_time\_$u$, life\_load\_$u$} & after the SOP ($u=0$) and each absorbed update \\
\cmd{t\_partition \dots\ t\_total} & CPU time of each phase and of the design point (s) \\
\bottomrule
\end{tabularx}
\end{center}

\subsection{Against other methods: simulated annealing and a genetic algorithm}
\label{sec:ex-baselines}

\begin{experimentbox}{baselines}{the heuristic against metaheuristics}
\exitem{Question.} Would a general-purpose search do better, given time?
\exitem{How.} \cmd{baselines results.csv trace.csv [-{}-sizes L] [-{}-seeds N]
[-{}-sa L] [-{}-ga L] [-{}-platform P] [-{}-no-local-search]} maps the
same applications (50, 100, 200, 300 chains, seeds 1 to 10, the default
generator) on the platform \cmd{default}, \cmd{slow} (zone cores at
25\,\%), \cmd{constraints} (replicas, zone-only chains, 8\,MB per zone
controller) or \cmd{domain}, with the four rules and with two
metaheuristics that solve the same allocation problem: the same
sub-chains, the same constraints, the same schedulability test, the same
first objective (the fewest cores).
\begin{itemize}
\item \textbf{Simulated annealing} (SA): from a first-fit mapping, moves
of one sub-chain or exchanges of two, accepted with the Metropolis rule;
budgets of $2\cdot10^4$, $2\cdot10^5$ and $10^6$ moves.
\item \textbf{Genetic algorithm} (GA): an individual is an order of the
sub-chains, decoded by first fit; tournament selection, order crossover,
swap mutation, elitism; 40 individuals; 5, 25 and 100 generations.
\end{itemize}
Both get the deadline local search of the heuristic on their best
mapping (unless \opt{no-local-search}), so that the chains on time compare
the searches. \cmd{include/mapping/search.h} gives every detail.
\exitem{Measures.} Per application and method: cores, lower bound,
schedulable cores, chains on time, CPU time; the trace file has every
improvement of the best mapping over time (the quality-against-time
figure of the paper).
\exitem{Reading.} Compare with the optimum of the exact solver for the
same applications (\cref{sec:ex-exact}). On the default platform the
rules take milliseconds, SA needs about $10^6$ moves (seconds) to match
them, and GA stays above them at 200 and 300 chains. On the domain
platform and with slow zones, GA (tens of seconds) finds fewer cores than
the rules; on the constrained platform no method comes close to the
exact solver at 200 chains.
\end{experimentbox}

\subsection{The exact solver}
\label{sec:ex-exact}

\begin{experimentbox}{exact}{the optimality gap}
\exitem{Question.} How far from the fewest possible cores is the heuristic?
\exitem{How.} The exact formulation (paper, section ``Exact formulation'')
finds the fewest cores for the sub-chains of the partitioning, with the
same constraints and schedulability test, as a constraint program
(OR-Tools CP-SAT) with lazy schedulability cuts.
\begin{shell}
build/bin/export\_instance 1 100 inst.json \quad\# seed, task-chains, file\\
python3 exact/exact.py inst.json -{}-time-limit 60 -{}-json result.json\\
python3 exact/gap.py build/bin/export\_instance -{}-sizes 20,50,100,200,300 -{}-seeds 10\\
python3 exact/gap.py build/bin/export\_instance -{}-sizes 20,50,100,200,300 -{}-seeds 10 \textbackslash\\
\quad -{}-configs domain,hybrid -{}-out results/exact\_architectures.csv
\end{shell}
\cmd{export\_instance} generates the application, partitions it once,
runs the four rules (with and without consolidation) and writes the
sub-chains, the candidate cores and the heuristic results.
\cmd{exact.py} prints the status (\cmd{OPTIMAL}, \cmd{FEASIBLE} at the
time limit, \cmd{INFEASIBLE}), the optimum and the time. \cmd{gap.py}
solves 150 instances on the zonal platform (default, slow zones,
constraints), checks every exact solution with the tool
(\cmd{export\_instance -{}-verify}) and writes \cmd{results/exact.csv};
the second run adds 100 instances on the domain and hybrid templates
(\cmd{export\_instance -{}-template T}). \opt{spare N} gives each ECU
$N$ candidate cores beyond the most any rule uses (2 by default); with 4,
no instance gets a smaller optimum.
\exitem{Reading.} The gap of a rule is its cores minus the optimum, in \%
of the optimum.
\end{experimentbox}

\subsection{Checking the bounds by simulation}

The analysis gives upper bounds. Two simulators check that no simulated
delay ever exceeds them.

\begin{experimentbox}{netsim}{the network, frame by frame}
\label{sec:ex-netsim}
\exitem{How.} \cmd{tests/netsim/netsim.py} simulates the frames of a flow
set at every output port: store-and-forward switches, CBS with the idle
slopes of the analysis, TAS gates, ATS regulators, FRER replication and
elimination, a saturated best-effort class, release jitter. It compares
every delay with its Network Calculus bound.
\begin{shell}
build/bin/net\_export map 0.3 1 flows.csv -{}-topology ring -{}-classes 3 -{}-frer\\
python3 tests/netsim/netsim.py flows.csv 20000 2 \quad\# horizon (\us), runs\\
sh tests/netsim/run\_netsim.sh build/bin/net\_export tests/netsim/netsim.py 20000 2
\end{shell}
\cmd{net\_export map <sigma> <seed>} exports the flows of a mapping,
\cmd{heavy <load> <seed>} random flows up to a load of every port;
\opt{config FILE} takes the platform and network of a configuration file.
The script runs 23 flow sets (star, tree, ring, custom; CBS, TAS, ATS;
FRER; the templates).
\exitem{Reading.} \cmd{max observed/bound} below 1 and
\cmd{violations = 0} on every line.
\end{experimentbox}

\begin{experimentbox}{syssim}{the whole system: cores, LET and network}
\label{sec:ex-syssim}
\exitem{How.} \cmd{tests/syssim/syssim.py} simulates a mapping written by
\opt{out} (\cmd{system.json}): every core under preemptive fixed
priorities with its LET task, every LET copy (implicit or strict LET, as
in the mapping), every message with a delay
up to its network bound, the cores on unsynchronized phases. It follows
each sensor sample through its chain to the first actuator job that uses
it, and compares the end-to-end delay and the maximum reaction time of
every chain with their bounds; half of the scenarios run every job at its
WCET and every message at its bound, the others draw them lower.
\begin{shell}
build/bin/mapping -{}-chains 200 -{}-out run \&\& \textbackslash\\
\quad python3 tests/syssim/syssim.py run/system.json 4\\
sh tests/syssim/experiment.sh build/bin/mapping results/syssim.csv 5 4
\end{shell}
The experiment script maps five applications of 200 chains on nine
settings (zonal star with FFDU and CADU, ring with FRER and three classes,
memory contention, tree with TAS, strict LET on the star and on the ring,
the domain and hybrid templates) and
writes every chain's bounds and observed values to \cmd{results/syssim.csv}.
\exitem{Reading.} \cmd{violations = 0} on every line; the ratio
observed/bound of each chain shows how tight the bounds are (1.000 for a
chain on one core at its WCETs: the bound is reached).
\end{experimentbox}
